\section{The Propers: Detailed Commentary}
\subsection{Introit: the Lamb slain and acclaimed}
The opening antiphon draws together Apocalypse 5:12 and 1:6. In the fifth chapter, John sees the Lamb standing as slain within the heavenly court. The Lamb is worthy to open the closed book because he was slain and redeemed a people. The heavenly worshippers acclaim him in the words selected for the Introit. Apocalypse 1 likewise names kingdom and priesthood in its account of Christ's faithful witness and saving blood.

The received commentary of Victorinus on Apocalypse 5 connects the opening of the testament with the Lamb's overcoming death and bringing humanity into life. It expounds the chapter as a whole, without explaining every term of the chant's selected verse. The ANF edition raises attribution and interpolation questions about this received commentary.\footnote{Received Victorinus, \emph{Commentary on the Apocalypse}, ch. 5, in \emph{Ante-Nicene Fathers} VII, with its editorial elucidations; see also ch. 1 on vv. 4--7.}

The Latin acclamation includes \emph{divinitatem}, divinity, among the things the Lamb is worthy to receive. The Epistle identifies him as the one through whom all creatures exist. Read together, these passages concern the acclaim of the divine Son, not his acquisition of a divine nature after death. The Introit's psalm verse asks for judgment and justice, so that the assembly's praise includes a petition for just rule.

\subsection{Collect: sin and divided peoples}
The Collect asks for the restoration of all things in the beloved Son, then names the peoples scattered by the wound of sin. The division for which it seeks healing is therefore both a relation to Christ and a relation among human families. It asks that they be gathered under his dominion and describes that dominion as gentle.

The Gospel describes allegiance to Christ as hearing his voice. Augustine explains how regeneration makes that hearing possible; Chrysostom distinguishes rescue from darkness from the further gift of transfer into the Son's kingdom. Their accounts help explain the change for which the Collect asks: people estranged by sin need to be made capable of life under Christ's rule.\footnote{Augustine, \emph{Tractates on John} 115.2,4; Chrysostom, \emph{Homilies on Colossians} II, on Col 1:12--15. These witnesses explain the biblical passages, not the composition of the Collect.}

\subsection{Epistle: rescue, creation and reconciliation}
The thanksgiving of Colossians 1:12--20 begins with the recipients. They have been made fit for an inheritance and transferred from darkness into the beloved Son's kingdom. It then identifies the Son as the one through whom all creatures exist. The person who saves them is therefore not one of the invisible creatures through whom they must climb to God. Visible and invisible creatures alike depend on him, including every heavenly power Paul names.

The passage uses \emph{firstborn} twice: of the Son's relation to creation and of his relation to the dead. The one through whom creatures exist is also the risen head of the Church. The latter part then names the Cross as the means of peace. Chrysostom emphasizes both the Creator's continuing preservation of creatures and the extraordinary condescension of his joining himself to the Church as its head.\footnote{Chrysostom, \emph{Homilies on Colossians} III, on Col 1:15--20, NPNF I.13.}

The next verses, Colossians 1:21--23, identify the recipients as former enemies and call them to perseverance. The Son's authority over all creation is distinct from the personal reconciliation they have received and must live. Universal authority does not imply automatic salvation irrespective of response. The second interpretation develops Chrysostom's answer to the phrase about things in heaven: what can their reconciliation mean if holy angels have not sinned?

\subsection{Gradual: the reach of justice}
Psalm 71:8 and 11 are two selected verses, not a continuous reading of verses 8--11. They describe dominion across seas and over every nation's rulers. The intervening and surrounding psalm describes homage and deliverance; verses 12--14 recount the king's rescue of the needy. The king who receives universal homage cares for the helpless.

Augustine begins his exposition with Christ as the true Solomon, the peacemaker whose kingdom exceeds that of the historical king. Bellarmine permits an immediate petition by David for Solomon's just government before unfolding its fulfillment in Christ. The difference affects how the opening petition is heard: one begins directly with the fulfillment, the other gives the historical royal prayer its own place within it. Both understand the reach and endurance of the psalm to surpass Solomon.\footnote{Augustine, \emph{Expositions of the Psalms}, Ps 71.1--2,13--16 (NPNF heading: Psalm 72); Bellarmine, \emph{Commentary on the Book of Psalms}, Ps 71, opening and comments 10--13 in the O'Sullivan delivery. Comment 10 treats biblical v. 11.}

\subsection{Alleluia: a kingdom that does not pass away}
Daniel 7 first displays beastly kingdoms, then the heavenly court and the human figure who receives dominion. The selected concluding clauses of verse 14 celebrate that dominion's endurance. The longer verse and the chapter explain whose dominion it is and how the saints belong to its future. The Mass does not sing the entire vision, but its unending kingdom has its meaning within that contrast between what rises from the earth and what God gives.

Jerome explains received power according to the human nature the eternal Son assumed and the humiliation described in Philippians 2. It does not imply an inferior divine nature. The same commentary denies that the saints' inheritance is exhausted by a temporary earthly victory. His account of their enduring kingdom helps explain the hope of shared heavenly reign in the Postcommunion.\footnote{Jerome, \emph{Commentary on Daniel}, VII, on 7:13--14 and 17--18; PL 25, cols. 670--671.}

\subsection{Gospel: the prisoner questions the judge}
The appointed scene is inside the governor's hall. Pilate asks whether Jesus is King of the Jews; Jesus asks where that description comes from. The exchange discloses the accusation and exposes the questioner's assumptions. Pilate is investigating a threat to authority. Jesus answers by identifying the source of his kingdom and the purpose of his coming into the world.

John 18 has already included Peter's sword and Jesus' refusal of it. In the trial, the absence of a fight to prevent his delivery becomes evidence for the kingdom's different origin. Jesus then speaks of testimony and hearing. The appointment ends there, at verse 37. Pilate's subsequent question about truth and his departure belong to the surrounding narrative, not to the appointed Gospel's final sentence. The assembly is left with a call to hear. The narrative identifies particular speakers and authorities; its conflicts do not assign guilt to Jewish people across generations.\footnote{John 18:10--11,28,33--40; 19:1--22; Augustine, \emph{Tractates on John} 115.1--5.}

\subsection{Offertory: intercession for the nations}
Psalm 2's promise of the nations follows rebellion against the Lord and his anointed. At verse 8 the Father commands the Son to ask and promises to grant the nations as his inheritance. Augustine relates this asking to the Incarnation, sacrifice and intercession. The Son's eternal dignity and his temporal work for humanity belong to one person. Understood through that intercession, the Offertory concerns the salvation of the nations the Son receives.\footnote{Augustine, \emph{Expositions of the Psalms}, Ps 2.7; Bellarmine, \emph{Commentary on Psalms}, Ps 2, comments on vv. 7--9.}

Verse 9's rod belongs to the psalm's context, not the Offertory's sung words. Bellarmine's account includes correction of sinful habits and final judgment. The passage demands accountability before the Son; it supplies no commission for worshippers to impose conversion by violence.

\subsection{Secret: the offered Son gives peace}
The Secret names human reconciliation and asks the Father that Christ, the Son offered in sacrifice, himself give unity and peace among the nations. The peoples whose division the Collect described are the recipients of this petition. Although the prayer is addressed to the Father, it identifies the Son as both the one offered and the giver of peace.

The Epistle identifies the blood of the Cross as the means of peace; the Secret asks for peace through the Son offered in sacrifice. Chrysostom's explanation of Colossians concerns precisely this means: reconciliation through Christ himself, his death and his Cross. On that understanding, the community asks to receive among the nations the peace his self-giving accomplishes.

\subsection{Preface: the royal priest}
The Preface names Christ eternal Priest and universal King and recounts his self-offering on the Cross. It describes the kingdom presented to the Father through truth and life, holiness and grace, justice, love and peace. These qualities concern both the life received from God and the relations of those who serve him. After recounting Christ's offering and kingdom, the assembly gives thanks with the heavenly host.

The Son whom the Epistle identifies as the sustainer of creation is here named as the priest who offers himself. The Preface recounts the subjection of creatures to his rule before the presentation of the kingdom to the Father. Pius XI's feast teaching likewise relates kingship to redemption and priesthood, while extending its moral claims to public life.\footnote{Pius XI, \emph{Quas primas} (11 December 1925), nos. 12--20.}

\subsection{Communion and Postcommunion: peace, service and shared reign}
The Communion takes selected clauses from Psalm 28:10--11. The psalm's storm, powerful divine voice and disrupted landscape precede its conclusion in royal permanence and peace. Augustine reads the Church as exposed to the world's flood while receiving peace in Christ. In his interpretation, the blessing can be received amid continuing adversity. Bellarmine preserves both natural and ecclesial readings of the psalm and relates present strength to final blessing.\footnote{Augustine, \emph{Expositions of the Psalms}, Ps 28.10--11 (NPNF heading: Psalm 29); Bellarmine, \emph{Commentary on Psalms}, Ps 28, conclusion.}

The Postcommunion names those already nourished by the sacrament who glory in serving under Christ's standard; it asks that they reign with him in heaven. Read beside the Gospel, its military image describes fidelity to the King whose servants do not fight to prevent his self-offering. Such service is also consistent with the peace requested at the altar. The nourishment already received can sustain their perseverance toward the lasting reign for which they pray.
